% !TEX root = ../diary-book.tex

\section{令和7年01月14日 シャワー測定法の更新可能頻度について}


今は 100 x 100 m2 の領域しか使っていないけれども、使えないかどうかはもう少し考える必要がある。

これまでのまとめを記述しておく。

到来宇宙線の頻度については PDG に載っているので、これに更新するとして、現状は wikipedia のものを使っている。
この頻度で陽子が 20 km 上空から入射していることを仮定している。

簡単のために $100 \times 100$~m$^2$の領域に検出器を 10~m 毎に全部で 100~台設置下とする。

フラックス $\phi$ を持つミュオンがある検出器で測定されたとき、